\documentclass[11pt]{article}

\usepackage[margin=0.5in]{geometry}
\usepackage{amsmath,amsfonts,amssymb,bm,mathrsfs,xspace}
\usepackage{xcolor}
\usepackage{enumitem}
\usepackage[normalem]{ulem}
\usepackage{amsthm}
\newtheorem{proposition}{Proposition}
\newtheorem{assumption}{Assumption}
\usepackage{cancel}

% Shared mathematical notation from paper/ex_shared.tex.
\newcommand{\RR}{\mathbb{R}}
\newcommand{\EE}{\mathbb{E}}
\newcommand{\PP}{\mathbb{P}}
\newcommand{\prox}{\mathrm{prox}}
\newcommand{\dist}{\operatorname{dist}}
\newcommand{\dom}{\mathrm{dom}}
\newcommand{\ran}{\mathrm{ran}}
\newcommand{\Id}{\mathrm{Id}}
\newcommand{\calB}{\mathcal{B}}
\newcommand{\calC}{\mathcal{C}}
\newcommand{\calD}{\mathcal{D}}
\newcommand{\calJ}{\mathcal{J}}
\newcommand{\calL}{\mathcal{L}}
\newcommand{\calO}{\mathcal{O}}
\newcommand{\calQ}{\mathcal{Q}}
\newcommand{\calS}{\mathcal{S}}
\newcommand{\bB}{\bm{B}}
\newcommand{\bD}{\bm{D}}
\newcommand{\bd}{\bm{d}}
\newcommand{\bG}{\bm{G}}
\newcommand{\bLmbd}{\bm{\Lambda}}
\newcommand{\bx}{\bm{x}}
\newcommand{\by}{\bm{y}}
\newcommand{\bz}{\bm{z}}
\newcommand{\be}{\bm{e}}
\newcommand{\bDelta}{\bm{\Delta}}
\newcommand{\bb}{\bm{b}}
\newcommand{\bh}{\bm{h}}
\newcommand{\ba}{\bm{a}}
\newcommand{\bg}{\bm{g}}
\newcommand{\bu}{\bm{u}}
\newcommand{\bC}{\bm{C}}
\newcommand{\bZ}{\bm{Z}}

\newcommand{\blue}[1]{{\ifmmode\color{blue}#1\else\normalfont\upshape\color{blue}#1\fi}}
\newenvironment{bluequote}{\begin{quote}\normalfont\upshape\color{blue}}{\end{quote}}
\newcommand{\red}[1]{\begin{color}{red}\itshape #1\end{color}}
\newcommand{\green}[1]{{\color{green!50!black}#1}}

\title{Response to Reviewers}
\author{}
\date{}

\begin{document}
\maketitle
\begingroup
\let\reviewitem\item
\renewcommand{\item}{\reviewitem\itshape}

For ease of reference, excerpts from the revised manuscript are shown in blue.

\vspace{1em}

\section{Reviewer 1}

\subsection{Reviewer's general assessment}

\noindent\textit{Summary.}
\red{The paper presents DualHash, a primal-dual reformulation for deep hashing with W-type quantization regularization. My understanding is that the main contribution is an adaptation of primal-dual splitting, combined with stochastic momentum/STORM updates, to deep hashing. The method uses the Fenchel conjugate of a specific W-type regularizer to obtain a closed-form dual proximal update. However, the convexity of the dual/conjugate side is standard.}

\subsection{Major issues}

\subsubsection*{(a) Relation between the primal-dual reformulation and the original objective}

\begin{itemize}\itshape
\item The key question is whether this transformed problem provides meaningful guarantees for the original nonconvex W-regularized hashing objective. I do not find this connection convincingly established. In particular, the paper should clarify whether there is strong duality, zero duality gap, or at least a rigorous stationarity relationship between the transformed primal-dual problem and the original objective.
\par\smallskip
\normalfont\upshape\textbf{Response:} 
\end{itemize}

\subsection{Minor issues}

\subsubsection*{Definition of the W-type regularizer and clarity of the setting}

\begin{itemize}\itshape
\item I also found the setting and writing unclear. The term ``W-type regularizer'' is used repeatedly but never properly defined; only one example is given. Even by the start of Section 3, the precise problem being studied is still not fully clear. If ``W-type'' denotes a class of regularizers, that class should be formally specified. If it refers only to the particular regularizer $h(z) = \lambda \lVert |z| - 1 \rVert_1$, the paper should state this much earlier and avoid presenting the results as applying to a broader family.
\par\smallskip
\normalfont\upshape\textbf{Response:} 
\end{itemize}

\subsubsection*{Length and organization}

\begin{itemize}\itshape
\item \red{Finally, the paper feels longer than necessary relative to its main contribution. Some of the space spent on standard preliminaries and routine descent arguments could be shortened, cited, or moved further into the supplement. This would allow the main text to focus more sharply on what is actually novel, especially the relationship between the primal-dual reformulation and the original nonconvex W-regularized hashing objective.}
\par\smallskip
\normalfont\upshape\textbf{Response:} 
\end{itemize}

\clearpage
\section{Reviewer 2}

\subsection{Reviewer's general assessment}

\noindent\textit{Opening summary.}
The aim of this article is to propose a new primal-dual stochastic approach, called DualHash, to address the challenge of `deep hashing' encoding within a non-convex optimisation framework. This has led to the development of two algorithms, DualHash-StoM and DualHash-Stom, based respectively on the principle of momentum acceleration and on a variance reduction technique. The authors present a complexity analysis of the two schemes before conducting experiments to compare their performances with existing methods on image datasets. The proposed article is relatively clear and well-referenced, with a comprehensive introduction. Its strongest and most relevant being the one dedicated to numerical experiments well detailled. However, the theoretical sections contain several weaknesses that need to be addressed in order to fully meet the journal's requirements. For all these reasons, \textbf{I return the article to the authors for major revisions.}
\medskip
\noindent\textit{1 Abstract/Introduction.}
Abstract and introduction appear generally well-written and well-cited to me . It is clear what the authors are trying to achieve in this paper.
\medskip
\noindent\textit{5 Complexity analysis.}
This section brings together a fairly large number of results, including intermediate lemmas and theorems, which may seem a bit exhaustive and even disjoined as it forces the reader to keep referring back and forth to the appendices. To such an extent, my comments relative to proofs postponed in Appendix C or D are covered in a dedicated section, at the end of this document. A few mathematical inaccuracies needs to be reworked and particularly the stochastic formalism of the paper needs to be a little refined. Finally, your complexity analysis involves a lot of parameters whose values tend to vary depending on the subsections, which can sometimes detract from the main point.
\medskip
\noindent\textit{9 Appendix C.}
The computational part of the proofs is generally carried out correctly, with only a few minor errors. However some justifications are missing or confused particularly, the way you pass to the conditional expectation.
\medskip
\noindent\textit{10 Appendix D.}
The strengths and weaknesses of the proofs contained in this section are similar to those in the previous appendix.

\subsection{Major issues}

\subsubsection{Stochastic formalism}

\begin{itemize}\itshape

\item (Assumption 5.1) In (iii), you use several notations that you forgot to introduce earlier, in particular the conditional expectation operator $\mathbb{E}_i$. To do so, you also need to clearly define a probability space and a filtration. That is why I would tend to add an extra paragraph and/or move here the one you wrote at the very beginning of the Appendix C.

\item Equalities or inequalities involving random variables (typically those containing conditional expectations) are `only' verified almost surely; \red{this is why it is important to specify this whenever necessary}.

\item (page 29) Notations $\mathbb{E}_{\xi^{[k]}}$ and $\mathbb{E}_k$ (rather written $\mathbb{E}_i$) introduced in the main paper seems to correspond to the same object. Also I would place this notation part in Section 5 instead (see my 3rd point of ``Complexity analysis section).\normalfont

\par\smallskip
\normalfont\upshape\textbf{Response:} We thank the reviewer for these careful comments regarding the stochastic formalism. We agree that in the original manuscript, the finite-sum average \(\mathbb{E}_i\) was not clearly distinguished from the conditional expectations \(\mathbb{E}_{\xi^{[k]}}\) and \(\mathbb{E}_k\). We apologize that our notation obscured this distinction. We have therefore rewritten the sampling assumptions explicitly in finite-sum form, unified the conditional-expectation notation as \(\mathbb{E}_k[\cdot]=\mathbb{E}[\cdot\mid\mathcal{F}_k]\), and stated the corresponding almost-sure convention. We have revised the beginning of Section~4.2, Section~5 and Appendices~C and~D accordingly, as detailed below.
\begin{enumerate}[label=\textbf{(\Roman*)},leftmargin=2.2em]

\item \textbf{Probabilistic framework and conditional expectation notation.} \red{In the submitted manuscript, Section~4.2 did not explicitly distinguish the sample index from the iteration index or define the mini-batch size in the sampling statement.} We have revised Section~4.2 accordingly \green{(Page 9, lines 262--265):}
\begin{bluequote}
Here and below, $i\in\{1,\ldots,n\}$ indexes training samples, whereas $k$ indexes algorithmic iterations. At iteration $k$, we sample a mini-batch $\mathcal J_k$ uniformly from all subsets of $\{1,\ldots,n\}$ of cardinality $b_k$, independently of the previous mini-batches [\ldots]
\end{bluequote}
We have replaced the notation \(\mathbb{E}_{\xi^{[k]}}\) with the unified notation \(\mathbb{E}_k\), and removed the ambiguous notation \(\mathbb{E}_i\) by writing the corresponding finite-sum averages explicitly. We have also moved the probabilistic framework to the beginning of Section~5, immediately before Assumption~5.1 (\green{Page 11, lines 1--10}):
\begin{bluequote}
Under the mini-batch sampling scheme introduced in Section~4.2, let $(\Omega,\mathcal F,\mathbb P)$ be the underlying probability space. For $k\geq1$, define $\mathcal F_k:=\sigma(\mathcal J_1,\ldots,\mathcal J_{k-1})$ with $\mathcal F_1$ trivial, as the filtration generated by the past mini-batches. We write $\mathbb E_k[\cdot]:=\mathbb E[\cdot\mid\mathcal F_k]$
for the conditional expectation given the past and $\mathbb E[\cdot]$ for the total expectation. All equalities and inequalities involving $\mathbb E_k[\cdot]$ are understood to hold almost surely, and we explicitly indicate ``a.s.'' whenever such conditional relations are stated.
\end{bluequote}
The redundant batch-history notation previously introduced in Appendix~C has been removed, and Appendices~C and~D now use the probabilistic notation defined in Section~5 
(\green{Page 31, lines 1--10}):
\begin{bluequote}
Before proceeding, throughout Appendices~C and~D, we adopt the notation, {the probabilistic setting,} and intermediate results from Section~5.
\end{bluequote}
\item \textbf{Finite-sum form of the sampling assumptions.} \red{The revised manuscript now writes the generic uniformly sampled finite dataset} \red{as
\[
\mathbb E_{\xi}[\phi(\bx;\xi)]
=\frac{1}{n}\sum_{i=1}^{n}\phi(\bx;\xi_i).
\]
Our problem has exactly this finite-sum structure, with
\[
F(\bx,\bB)=\frac{1}{n}\sum_{i=1}^{n}F_i(\bx,\bb_i).
\]}
Accordingly, we have rewritten Assumption~5.1(iii) directly in finite-sum form \green{(Page 12, lines 1--5):}
\begin{bluequote}
    For any $\bB \in \RR^{n\times d}$, there exists a constant $\sigma^2 \geq 0$ such that
\[
\frac{1}{n}\sum_{i=1}^{n}\nabla_{\bx}F_i(\bx,\bb_i)
=\nabla_{\bx}F(\bx,\bB),
\qquad
\frac{1}{n}\sum_{i=1}^{n}
\left\|\nabla_{\bx}F_i(\bx,\bb_i)
-\nabla_{\bx}F(\bx,\bB)\right\|^2
\leq\sigma^2. 
\]
\end{bluequote}
We have also revised Assumption 5.2 on mean-squared smoothness \green{(Page 16, lines 6--10):}
\begin{bluequote}
There exists $L_{F} > 0$ such that for all $\bx, \bx' \in \RR^{q}$ and $\bB, \bB' \in \RR^{n\times d}$, it holds that
\begin{align}
\frac{1}{n}\sum_{i=1}^{n}
\left\| \nabla_{\bx} F_{i}(\bx,\bb_i)
- \nabla_{\bx} F_{i}(\bx',\bb_i')\right\|^{2}
\leq L^{2}_F \left[ \|\bx - \bx'\|^{2}
+ \|\bB - \bB'\|^{2} \right]. \notag
\end{align}
\end{bluequote}

\item \textbf{Conditional expectations and almost-sure statements in the analysis.} 
\red{We agree xxx.} \green{(Page 13, lines 1--12)}, 
Please see our reponse to \red{xx}.
\end{enumerate}

\item {\itshape (page 13) Give more details to justify
$\mathbb{E}[\|\be^k\|^2]\le\sigma^2/b_k$.}

\par\smallskip
\normalfont\upshape\textbf{Response:}
We thank the reviewer for this comment to make our presentation rigiously. With the revised probabilistic
framework and conditional-expectation notation above for our algorithms, we give the detailed derivation, which has been added to \green{ Appendix~C.1  (Page XX, lines XX--XX)}:
\begin{bluequote}
Recall that the stochastic gradient error $\be^k$ of Algorithm 4.1 is defined as
\[
\be^k
=
\frac{1}{b_k}\sum_{i\in\calJ_k}
\left(
\nabla_{\bx}F_i(\bz^k,\bb_i^k)
-
\nabla_{\bx}F(\bz^k,\bB^k)
\right).
\]
Since $(\bz^k,\bB^k)$ is $\mathcal F_k$-measurable and
$\calJ_k$ is sampled uniformly without replacement and independently
of $\mathcal F_k$, define
\[
\bm{v}_i^k
:=
\nabla_{\bx}F_i(\bz^k,\bb_i^k)
-
\nabla_{\bx}F(\bz^k,\bB^k).
\]
Conditionally on $\mathcal F_k$, Assumption~5.1(iii)
gives
\[
\frac{1}{n}\sum_{i=1}^n \bm{v}_i^k=\bm{0},
\qquad
\frac{1}{n}\sum_{i=1}^n\|\bm{v}_i^k\|^2\leq\sigma^2
\quad\text{a.s.}
\]
For $n\geq2$, the variance formula for uniform sampling without replacement
therefore yields
\[
\mathbb{E}_k[\be^k]=\bm{0},
\]
and
\[
\mathbb{E}_k\!\left[\|\be^k\|^2\right]
=
\frac{n-b_k}{n-1}\frac{1}{b_k}
\frac{1}{n}\sum_{i=1}^n\|\bm{v}_i^k\|^2
\leq
\frac{\sigma^2}{b_k}
\quad\text{a.s.}
\]
The case $n=1$ is immediate since $\be^k=\bm{0}$. Taking the full
expectation and using the tower property gives (5.11).
\end{bluequote}

We have also summarized this argument in the main text and added an equation label for \green{later reference (5.11)(Page 14, lines XX--XX)}.
\end{itemize}

\subsubsection{Boundedness and lower bounds}

\begin{itemize}\itshape
\item (Lemma 5.1/proof) Having $J$ coercive and lower-bounded is generally not enough to directly obtain the boundedness of $(\bx^k,\bB^k)$. You classically also need to prove that sequence $\{J(\bx^k,\bB^k)\}$ converges (or at least its boundedness). Note also that having boundedness of all level-sets is equivalent to coercivity. Finally, the coercivity of $J$ is not assumed in your assumption 5.1.
\par\smallskip
\normalfont\upshape\textbf{Response:} We thank the reviewer for the important comment about the boundedness argument. While we agree that deducing boundedness of the iterates in Lemma~5.1 from coercivity of the objective function $J(\bx,\bB)$ requires a uniform upper bound on the objective values along the iterates, our complexity analysis actually does not require the iterates themselves to be bounded. Rather, it requires lower bounds on the potentials used in the analysis. This was the original motivation for introducing Lemma~5.1.

Regarding coercivity, we already stated in the submitted manuscript (p.~7) that the objective $J$ in~(3.1),
\[
J(\bx,\bB)=
\underbrace{\frac1n\sum_{i=1}^n\ell_{s,i}(\bx)}_{=:\,\bar{\ell}_s(\bx)}
+\frac{\mu}{2}\|\bx\|^2
+\frac{\gamma}{2n}\sum_{i=1}^n\|\mathcal D_i(\bx)-\bb_i\|^2
+h(\bB),
\]
is coercive. However, to ensure this property through the model structure, we impose the lower-boundedness assumption on the average loss rather than directly on the objective:
\[
\underline{\ell}:=\inf_{\bx}\bar{\ell}_s(\bx)>-\infty.
\]
Together with the model structure, this condition ensures both coercivity and lower-boundedness of $J$. We therefore refine Assumption~5.1(i) accordingly. This refinement rules out cancellation of the growth of the regularization terms by the loss and justifies the coercivity stated in the manuscript.

Under the refined Assumption~5.1(i) and the appropriate parameter conditions, the estimates in our analysis yield the required lower bounds on the potentials. We have made these estimates explicit in Proposition~\ref{prop:lyapunov-lower-bound} and its proof in the ``Lower bounds for the Lyapunov functions'' section below. These adjustments preserve the stated complexity orders and the main theoretical contributions of the manuscript.
\item (Lemma 5.7) Why $R\in 2,...,T+1$ instead of $R\in 1,...,T$ ? Furthermore, I am not particularly in favour of using uniform sampling to present a rate; I would prefer to use the quadratic metric $\frac{1}{T}\sum_{k=1}^{T}\mathbb{E}[dist^2(...)]$. Finally, to use $\Psi_m^*$ you need at least to justify the lower-boundedness of $\{\Psi_m(\mathscr{Q}^k)\}$ which doesn't seem straightforward.
\par I would first introduce $L_\Delta$ before mentioning $S_1,S_2,S_3,S_4$ (these constant actually don't appear in (5.17) as such).
\par\smallskip
\normalfont\upshape\textbf{Response:} We thank the reviewer for these remarks. Regarding the index set and the metric: the shift $R\in\{2,\ldots,T+1\}$ was unnecessary, since the averaged bound in the proof runs over $k=1,\ldots,T$. In the revised manuscript, $R$ is selected uniformly from $\{1,\ldots,T\}$, independently of the iterates, in Lemma 5.7, Theorem 5.2, and the output step of Algorithm 4.1; the same correction is made in Lemma 5.10, Theorem 5.3, and Algorithm 4.2. Following the reviewer's preference, Lemma 5.7 now also states the bound for the quadratic metric $\frac{1}{T}\sum_{k=1}^{T}\mathbb{E}[\operatorname{dist}^2(0,\partial\calL(\bx^k,\bB^k,\bLmbd^k))]$, which equals $\mathbb{E}[\operatorname{dist}^2(0,\partial\calL(\bx^R,\bB^R,\bLmbd^R))]$ under this choice of $R$. Regarding $\Psi_m^*$, the lower bound of $\{\Psi_m(\calQ^k)\}$ is established in Proposition~\ref{prop:lyapunov-lower-bound}(i); it uses only the global $L_F$-smoothness, $F\ge\underline{\ell}$, the box constraint on $\bLmbd^k$, and the output condition of Appendix~B, and it does not require the iterates to be bounded. We have also moved the definition of $L_\Delta$ to the front of the ``where'' clause of Lemma~5.7, before $S_1,\ldots,S_4$ are mentioned.
\item (Lemma 5.10) Similar comment as for Lemma 5.7, especially the fact you need to justify the lower-boundedness of $\{\Psi_r(\mathscr{Q}^k)\}$ to be able to invoke $\Psi_r^*$.
\par\smallskip
\normalfont\upshape\textbf{Response:} We thank the reviewer. As for Lemma~5.7, the lower bound needed to define $\Psi_r^*$ is given in Proposition~\ref{prop:lyapunov-lower-bound}(ii). It bounds from below the augmented potential $W_k=\mathbb{E}[\Psi_r(\calQ^k)]+\frac{\eta_T}{\rho_T}\mathbb{E}\|\be^k\|^2$, which contains the estimation-error term that the descent inequality produces, and $\Psi_r^*$ is defined from this bound.
\end{itemize}

\subsubsection{Errors in Appendix C (DualHash-StoM)}

\begin{itemize}\itshape
\item (Lemma 5.5/proof) See my comments relative to Appendix C.
\item (Lemma 5.6/proof) See my comments relative to Appendix C.
\item (Lemma 5.7/proof) See my comments relative to Appendix C.
\item (Theorem 5.2/proof) See my comments relative to Appendix C.
\par\smallskip
\normalfont\upshape\textbf{Response:} These four comments refer to the Appendix~C comments listed below, which are answered individually: Lemmas~C.1--C.3, the proofs of Lemmas~5.5--5.7, and the proof of Theorem~5.2.
\item (Lemma 5.6/proof) In (C.7) you don't pass to conditional expectation correctly ; this should be be $\mathbb{E}_{\xi^{[k+1]}}$ before every terms. Fortunately, this doesn't seem compromise your proof to the extent that the ``wrong'' terms in $\mathbb{E}_{\xi^{[k]}}$ in your inequality seems measurable (to check) and can therefore be written without a conditional expectation, for example $\cancel{\mathbb{E}_{\xi^{[k]}}[\Delta_{\bx}^k]}\to\mathbb{E}_{\xi^{[k+1]}}[\Delta_{\bx}^k]=\Delta_{\bx}^k$. Moreover, remind the reader you also use $\mathbb{E}_{\xi^{[k+1]}}[\|e_k\|^2]\le\sigma^2/b_k$ (to justify).
\par Finally, to reach the inequality after (C.7), you need to invoke Lemma C.2. For this reason, I would place Lemma C.2 before Lemma C.1.
\par\smallskip
\normalfont\upshape\textbf{Response:} We thank the reviewer. In the proof of Lemma~5.6, we now pass to the total expectation $\EE[\cdot]$ right after the pathwise Lyapunov inequality, in line with the reviewer's own suggestion for Lemma~D.1 to work with in-expectation statements. No conditional expectation remains in this proof, so the measurability of the terms in $\EE_{\bm{\xi}^{[k]}}$ no longer needs to be checked. The bound $\EE[\|\be^k\|^2]\le\sigma^2/b_k$ used there is justified in our responses on the probabilistic setting and to the comment on page~13 above. The proof now invokes Lemma~C.2 explicitly where the coefficients are substituted. We keep the order of the lemmas, since swapping them would renumber Lemma~C.1 and Lemma~C.2, which are cited throughout the reviewer's comments and this letter.
\item (Lemma C.2/proof) Before the calculus relative to $K_5-K_3-C_3-C_4$, remind the reader equalities $K_5=\frac{1-\alpha_k}{\eta_k}-L_F$, $C_3=K_4$, $\overline{L}_F=2L_F+K_3+K_4$. In the calculus of $C_4-K_6$, the 3rd line is even an equality. Moreover, this should be $\beta_k$ instead of $\beta$ everywhere.
\par Regarding the last part, delete the two ``$\ge 0$'' after $C_B$ and $\frac{2}{L_F}C_B$ expression respectively. Otherwise, this would be a logical error as, at this stage of the proof, you have not yet shown that $C_B$ is positive. Moreover, I found $\tau\in\left(0,6\left(\frac{\sqrt{56}-2}{13}\right)\delta\frac{1}{L_F}\right)$. Finally, could you provide additionnal details to justify why $C_B>0$ ensures $C_2>0$ ?
\par\smallskip
\normalfont\upshape\textbf{Response:} We thank the reviewer for the careful verification. In the revised proof of Lemma C.2, we recall $K_5=\frac{1-\alpha_k}{\eta_k}-2L_F$, $C_3=K_4$, and $\overline{L}_F=2L_F+K_3+K_4$ before the computation of $K_5-K_3-C_3-C_4$; we state $\alpha_k=\alpha$ and $\beta_k=\beta$ from the parameter conditions before computing $C_4-K_6$, which makes the use of the constant $\beta$ consistent with the definition of $C_4$; and the third line of that computation is now written as an equality. We have also removed the two premature assertions ``$\ge 0$''.

Regarding the range of $\tau$: with $\delta=1/6$ and $t=\tau L_F$, the general expression of $C_{\bB}$ gives
\[
\frac{2}{L_F}C_{\bB}=\frac{1}{t}-\frac{13}{2}t-2=\frac{1}{t}\Big(1-2t-\frac{13}{2}t^2\Big).
\]
The first term was misprinted as $\frac{2}{t}$ in the submitted manuscript. The value $\frac{\sqrt{56}-2}{13}$ found by the reviewer is exactly what the misprinted coefficient produces, whereas the corrected coefficient gives $C_{\bB}>0$ if and only if $0<\tau<\frac{\sqrt{30}-2}{13L_F}$, which is the range stated in the manuscript. We have corrected the coefficient and written the range as an open interval. We have also added the requested justification: $C_1=K_1=3\delta\tau^{-1}>0$ and $C_2=\frac{1}{2}(K_2-K_1+K_3)=C_{\bB}+K_3>0$, because $K_3=\tau L_F^2/\delta>0$. Finally, we state explicitly that $\tilde\delta=c_\delta\delta^2$ with $0<c_\delta<36\big(\frac{\sqrt{30}-2}{13}\big)^2$, so that $\tau\le\sqrt{\tilde\delta}/L_F$ implies the above range.
\item (Lemma C.3) In (C.10), I think the correct expression is $S_4^k=6(L_F^2\beta_k^2+\underline{2}\eta_k^{-1}\alpha_k^2)$. To maintain consistency in the structure of your terms, I would rewrite $S_1=4(\tau^{-2}+3)$.
\par\smallskip
\normalfont\upshape\textbf{Response:} We thank the reviewer. The coefficient $6$ in front of $\eta_k^{-2}\alpha_k^2$ was indeed a typo: the proof (estimate of $\|\bd_{\bx}^k\|^2$) gives $6\eta_k^{-2}\|\bx^k-\by^k\|^2=12\eta_k^{-2}\alpha_k^2\Delta_{\bx}^k$, so that $S_4^k=6(L_F^2\beta_k^2+2\eta_k^{-2}\alpha_k^2)$, which we have corrected. We note that the exponent of $\eta_k$ is $-2$, since the term arises from $\eta_k^{-2}\|\bx^k-\by^k\|^2$. We have also rewritten $S_1=4(\tau^{-2}+3)$ as suggested.
\item (Lemma 5.7/proof) Remind the reader $\eta_k$ is constant to justify $S_3^k=S_3$, $S_4^k=S_4$. Since all the terms are positive, I think it would be more appropriate to bound (C.9) using $\max\{S_1,S_2,S_3,S_4\}$ instead of $\max\{S_1,S_2\}+\max\{S_3,S_4\}$.
\par When you sum over $k=1,...,T$, term $(3+L_\Delta)\frac{\sigma^2}{b_k}$ is incorrect and should be $\frac{(3+L_\Delta)\sigma^2}{L_FT}\sum_{k=1}^{T}\frac{1}{b_k}$.
\par In (5.17), in the upper bound, you thus should obtain $\frac{C_\sigma\sigma^2}{T}\sum_{k=1}^{T}\frac{1}{b_k}$ with $C_\sigma=\frac{3+L_\Delta}{L_F}$.
\par\smallskip
\normalfont\upshape\textbf{Response:} We thank the reviewer for identifying this error and agree that, for varying batch sizes, the averaged bound must retain $\frac{1}{T}\sum_{k=1}^{T}\frac{1}{b_k}$. Tracing this term back also revealed an inconsistency in Lemma C.3: its proof yields the noise term $3\|\be^k\|^2$, whose expectation is bounded by $3\sigma^2/b_k$, so the factor $1/L_F$ in the statement was a typo. We have removed it. Combining the corrected Lemma C.3 with the Lyapunov descent inequality, whose noise term $\sigma^2/(b_kL_F)$ is unchanged, gives the corrected bound (5.17):
\[
\frac{1}{T}\sum_{k=1}^{T}\mathbb{E}\big[\operatorname{dist}^2(0,\partial\calL(\bx^k,\bB^k,\bLmbd^k))\big]\le\frac{L_\Delta\Delta_1}{T}+\frac{C_\sigma\sigma^2}{T}\sum_{k=1}^{T}\frac{1}{b_k},\qquad C_\sigma=3+\frac{L_\Delta}{L_F}.
\]
The constant differs from the reviewer's $(3+L_\Delta)/L_F$ only because of the corrected Lemma C.3; the reviewer's expression is what the uncorrected statement gives. The corresponding steps in the proof of Lemma 5.7 have been revised.
\item (Theorem 5.2/proof) The previous mistake propagates into the proof of Theorem 5.2.
\par\smallskip
\normalfont\upshape\textbf{Response:} We thank the reviewer. With the constant batch size $b_k=b=c_bT$ used in Theorem 5.2, we have $\frac{1}{T}\sum_{k=1}^{T}\frac{1}{b_k}=\frac{1}{c_bT}$, so the corrected (5.17) gives $\mathbb{E}[\operatorname{dist}^2(0,\partial\calL(\bx^R,\bB^R,\bLmbd^R))]\le\frac{1}{T}\big(L_\Delta\Delta_1+C_\sigma\frac{\sigma^2}{c_b}\big)$. This is the bound used in the proof of Theorem 5.2, which we have updated to state this substitution explicitly. Hence the iteration bound and the $\calO(\varepsilon^{-4})$ SFO complexity in Theorem 5.2 remain unchanged.
\end{itemize}

\subsubsection{Errors in Appendix D (DualHash-StoRM)}

\begin{itemize}\itshape
\item ((Lemma 5.10)/Theorem 5.3 proofs) See my comments relative to Appendix D.
\par\smallskip
\normalfont\upshape\textbf{Response:} This comment refers to the Appendix~D comments listed below, which are answered individually: Lemmas~5.9, D.1, D.2, and D.4, and the proofs of Lemma~5.10 and Theorem~5.3. The two notation comments of Appendix~D, on $|\mathcal{I}_{k+1}|$ and on $C_\delta$, are answered under Notation and terminology.
\item (Lemma 5.9/proof) I don't really catch the interest of intermediate variable $s$ here. You should write (D.1) directly. Moreover, (D.1) is a simple direct consequence of Lemma 5.4 combined with Lemma 5.8. Why make things more complicated with intermediate equations (5.9), 5.7) ?
\par In (D.2), I don't understand where the term $(C_4-K_6)\Delta_x^k$ comes from.
\item (Lemma D.1) In Inequality of Lemma (D.1) and after a quick check of the proof of Lemma 8 of [49], I think this should be $\mathbb{E}_{\xi^{[k+1]}}[\|\be^k\|^2]$ instead of $\mathbb{E}_{\xi^{[k]}}[\|\be^k\|^2]$ (similar confusion as in the proof of (C.7)) . If it is so, you can hardly consider $\mathbb{E}_{\xi^{[k]}}[\|\be^k\|^2]$ as metric for the rest of your analysis. To avoid, this kind of problem, I think you should work with $\mathbb{E}[\|\be^k\|^2]$ metric instead and present an ``in-expectation'' version of Lemma D.1 i.e. of the form of :
\[
\mathbb{E}[\|\be^{k+1}\|^2]\le ...\mathbb{E}[\|\be^k\|^2]+...+...\mathbb{E}[\Delta_{\bx}^{k+1}].
\]
The resulting proof should therefore not undergo any significant changes. That is why, in the remaining proofs/results to review, I will assume that the inequalities are passed to the expectation rather than conditional expectation.

    \blue{
    \begin{align}
        \EE[\| \be^{k+1}\|^{2}] \leq \left(1 - \rho_{k}\right)^{2}\EE[\| \be^{k}\|^{2}] + \frac{2 \rho_{k}^{2}\sigma^{2}}{b_{k+1}}+ \frac{4(1 - \rho_{k})^{2}L_{F}^{2}}{b_{k+1}}\EE\left[\Delta_{\bx}^{k+1}+\blue{\Delta_{\bB}^{k+1}}\right]. \notag
    \end{align}}
\item (Lemma D.2/proof) For the second inequality of (D.7), the multiplicative factor in front of the first term should remain $\frac{1}{8(1-\rho_k)\eta_1L_f^2}$ instead of $\frac{1}{8\eta_1L_f^2}$, unless you are able to prove that $\mathbb{E}[\|\be^k\|^2]-\mathbb{E}[\|\be^{k+1}\|^2]+\frac{2\rho_k^2\sigma^2}{|\mathscr{I}_{k+1}|}\le 0$. This propagates into the rest of the proof.
\par Using (D.3) and more specifically $\rho_k\ge 8\eta_1\eta_kL_F^2$, I obtain $\frac{\rho_k}{8\eta_1L_F^2}-\frac{\eta_k}{2}\ge\frac{\eta_k}{2}$.
\par As in the proof of Lemma 5.7, you here need to justify that $\{\mathbb{E}[\Psi_r(\mathscr{Q}^k)]\}$ is non-increasing.
\par\smallskip
\normalfont\upshape\textbf{Response:} We thank the reviewer for this careful check. The reviewer is right that the factor $\frac{1}{8(1-\rho_k)\eta_1L_F^2}$ cannot simply be replaced by $\frac{1}{8\eta_1L_F^2}$, since the bracket it multiplies need not be nonpositive. We now apply the bound $1-\rho_k\le 1$ one step earlier, to the coefficient $4(1-\rho_k)L_F^2$ of $\mathbb{E}[\Delta_{\bx}^{k+1}]$ in the recursion of Lemma D.1, before rearranging. This gives $-\mathbb{E}[\Delta_{\bx}^{k+1}]\le\frac{1}{4L_F^2}\big((1-\rho_k)\mathbb{E}[\|\be^k\|^2]-\mathbb{E}[\|\be^{k+1}\|^2]+\frac{2\rho_k^2\sigma^2}{b_{k+1}}\big)$ and therefore yields the second inequality of (D.7) directly with the factor $\frac{1}{8\eta_1L_F^2}$, so the rest of the proof is unaffected. Regarding (D.3): both the reviewer's bound $\frac{\rho_k}{8\eta_1L_F^2}-\frac{\eta_k}{2}\ge\frac{\eta_k}{2}$ and the bound $\frac{\rho_k}{8\eta_1L_F^2}-\frac{\eta_k}{2}\ge\frac{\rho_k}{16\eta_1L_F^2}$ used in the manuscript follow from $\rho_k\ge 8\eta_1\eta_kL_F^2$; we keep the latter because the subsequent summation requires the weight $\rho_k$. While revising this proof, we also restored the term with $k=1$ that had been omitted in the final summation, so that the last sum in (D.4) now runs over $k=1,\ldots,T$; this adds a term of the same order and does not change the rate in Theorem 5.3.
\item (Lemma D.4 proof) I would justify the first term in (D.12) in a similar way to what you did for the second one, this time by invoking the second inequality in (5.23) and the fact that $\widetilde{S}_4\ge S_1,S_2$. Moreover, I think this is $\widetilde{S}_4$ everywhere instead of $S_4$.
\par It appears more logical me to substitute (D.13) into (D.12) before passing to expectation.
\par In (D.15), I think this is $\frac{2(C_\delta^1+4)}{\eta_1}$ instead of $\frac{2(C_\delta^1+4)}{S_4\eta_1}$.
\par The end of the proof as well as those of Lemma 5.10, Theorem 5.3 are based on the potential new (D.4) inequality. To be reworked accordingly.
\end{itemize}

\subsubsection{ Problem formulation and algorithm}

\begin{itemize}\itshape
\item (page 7) If I have clearly understood, problem (P) is a regularized version of (3.1) which is supposed to better leverage the structure of the initial problem. However, there seems to be a confusion relative to some variables definitions otherwise I don't really understand the point of introducing $P(\bx,\bB)$, (since $\bm{b}_i=\bm{D}_i(\bx)$, we have $\|\bm{b}_i-\bm{D}_i(\bx)\|^2=0$ and therefore $P(\bx,\bB)=0$). Finally, could you provide further details which would explain the interest of term $P(\bx,\bB)$ (it looks like a coupling one at first glance), don't hesitate to use references if necessary.
\item (page 9) In (4.5) first line, to avoid any confusion I would write ``$(\tau^{-1}h^*)$'' instead of ``$(1/\tau h^*)$''. Moreover, I think a $1/\tau$ factor is missing before $(2\bB^{k+1}-\bB^k)$.
\item (Algorithms 4.1,4.2) Reading through your pseudo-code, one might think the choice of every batch size is determined before the algorithm is run. In the case of a constant batch size $b_k=b$, that makes perfect sense otherwise, it seems a bit strange to me.
\end{itemize}

\subsubsection{Claims and applicability}

\begin{itemize}\itshape
\item (page 4) Regarding to Table 1.1, your paper seems to be the first one to conduct a complexity analysis for deep hashing methods. Are you totally sure of this ?
\item Regarding your Theorem 5.2, having a good $\epsilon$-precision thus requires to have a number of samples $b$ large enough (as $b=c_bT$) under the constraint to be smaller than $n$. Is it really achievable and not too restrictive in practice ?
\end{itemize}

\subsection{Minor issues}

\subsubsection{Proof exposition and organization}

Some important equations and formulations need to be referenced; the detailed responses follow.
\begin{itemize}\itshape
\item (page 5) I would give a number to the equation relative the reformulation of Problem (1.1). Also explain what the functions $f$ and $c$ are supposed to represent (for example is $f=1/n\sum f_i$ ? ...).
\par\smallskip
\normalfont\upshape\textbf{Response:}
The submitted manuscript did not number the deep-hashing formulation or identify the roles of $f$ and the network outputs clearly.
We have revised Section~3 accordingly \green{(Page 7, lines 191--204):}
\begin{bluequote}
Deep hashing aims to learn network parameters $\bx \in \RR^{q}$ that generate binary codes preserving semantic similarity in the Hamming space. A direct formulation is
\begin{align}
\min_{\bx \in \RR^{q}}\ \frac{1}{n}\sum_{i=1}^{n}\ell_{s,i}\!\big(\operatorname{sgn}(\bC_i(\bx))\big).
\end{align}
where $f_i(\bx)=\ell_{s,i}(\bx)+\frac{\mu}{2}\|\bx\|^2$ with $\mu>0$ to ensure coercivity. Thus, $f(\bx):=n^{-1}\sum_{i=1}^{n}f_i(\bx)$ denotes the averaged smooth loss, while $\bD_i(\bx)$ denotes the continuous network output for sample $i$.
\end{bluequote}

\item (page 6) Remind the reader that $\mathbf{d}^k$ correspond to the current direction at iteration $k$ (you might even recall the generic scheme $\bx^{k+1}=\bx^k-\eta_k\bd^k$). Moreover, give a number to the equation relative to $\bd^k$ update's rule for StoRM algorithm.
\par\smallskip
\normalfont\upshape\textbf{Response:}
The submitted manuscript did not distinguish the stochastic update direction from the subgradient used in the stationarity analysis or provide a numbered StoRM recursion.
We have clarified the direction notation and added the StoRM recursion in Sections~2 and~4.2 \green{(Page 6, lines 177--184; Page 9, lines 241--245):}
\begin{bluequote}
A standard momentum-based stochastic first-order method constructs the stochastic direction
\[
\bg^{k}=\rho_{k-1}\nabla\phi(\bx^{k},\xi^{k})+(1-\rho_{k-1})\bg^{k-1},\qquad \rho_{k-1}\in[0,1].
\]
At iteration $k$, $\bd^k$ denotes the stochastic direction used in the generic update $\bx^{k+1}=\bx^k-\eta_k\bd^k$.
\end{bluequote}

\item (page 6) Give a number to the equation relative to deep hashing + typo under the min `` $x\underline{\in}\mathbb{R}^q$ ''.
\par\smallskip
\normalfont\upshape\textbf{Response:}
The submitted manuscript left the direct deep-hashing problem unnumbered and contained the typo $x\underline{\in}\mathbb{R}^{q}$.
We have numbered the formulation and restored the membership symbol in Section~3 \green{(Page 7, lines 191--204):}
\begin{bluequote}
\begin{align}
\min_{\bx \in \RR^{q}}\ \frac{1}{n}\sum_{i=1}^{n}\ell_{s,i}\!\big(\operatorname{sgn}(\bC_i(\bx))\big).
\end{align}
\end{bluequote}

\item (Relation (C.8)) Remind the reader the expression of $\bd^k$ (this can be throughout a reference to an equation). Regarding your definitions relative to StoM, do we have $\bG^k=-\bd^k$ ? If not, how do you define $\bd^k$ otherwise (in section 2, you only remind the update's rule relative to StoRM) ? Furthermore, it is not clear how you define $\bd_{\bx}^k,\bd_{\bB}^k,\bd_{\bLmbd}^k$.
\par\smallskip
\normalfont\upshape\textbf{Response:}
The submitted manuscript used $\bd^k$ both for an update direction and for a Lagrangian subgradient, leaving the components in (C.8) unclear.
We have separated these roles and defined the subgradient components explicitly in Appendix~C \green{(Page 32, lines 951--956):}
\begin{bluequote}
The vector $\bd^{k}$ below is a subgradient of $\calL$ used only to measure stationarity; the update directions of the algorithms are $\bg_m^k$ and $\bg_s^k$. For $\bd^{k}\in\partial\calL(\bx^{k},\bB^{k},\bm{\Lambda}^{k})$, let us denote $\bd^{k}=(\bd_{\bx}^{k},\bd_{\bB}^{k},\bd_{\bm{\Lambda}}^{k})$ and
\begin{align}
\label{Eq: the components of the subgradient}
\begin{cases}
\bd_{\bx}^{k}=\nabla_{\bx}F(\bx^{k},\bB^{k}),\\
\bd_{\bB}^{k}=\nabla_{\bB}F(\bx^{k},\bB^{k})+\bm{\Lambda}^{k},\\
\bd_{\bm{\Lambda}}^{k}=\bB^{k}-\mathcal{G}_{\bm{\Lambda}}^{k},
\end{cases}
\end{align}
\end{bluequote}
\item (Lemma 5.2/proof) Some objects as $G$ function is not explicitly defined, making your proof a bit hard to follow. Moreover, which subdifferential property do you use to conclude ? To rework.
\par\smallskip
\normalfont\upshape\textbf{Response:}
The submitted proof did not make the coercivity and criticality arguments explicit.
We have expanded the proof of Lemma~5.2 to state the lower bound, dual-box boundedness, and subdifferential conclusion \green{(Page 12, lines 321--334):}
\begin{bluequote}
Since $f_i=\ell_{s,i}+\frac{\mu}{2}\|\cdot\|^2$ with $\mu>0$, we have $n^{-1}\sum_i f_i(\bx)\geq\underline{\ell}+\frac{\mu}{2}\|\bx\|^2$. The penalty term $P(\bx,\bB)$ is nonnegative, and $h(\bB)=\lambda\sum_{i,j}\bigl||B_{ij}|-1\bigr|\geq\lambda\|\bB\|_1-\lambda nd$.
\end{bluequote}
\item (Lemma 5.4/proof) I don't really understand how you reach inequality (5.9). Provide far more details (intermediate inequalities, mention descent lemma whenever it is used, how you deal with the conjugate quantities...). Assuming (5.9) holds, the rest of the proof is correct. Moreover, why did you choose to include this proof here while those of its counterpart (iteration progress on $\bx$) is presented in the appendix ?
\par\smallskip
\normalfont\upshape\textbf{Response:}
The submitted proof skipped the smoothness substitution and the treatment of the stochastic-direction error in the one-step estimate.
We have expanded the proof and moved it with the analogous proof to Appendix~C \green{(Page 33, lines 969--1000):}
\begin{bluequote}
Then, applying the auxiliary lemma to $F(\cdot,\bB^{k})$, which is $L_F$-smooth, with $\bx=\bx^{k}$, $\by=\by^{k}$, $\eta=\eta_{k}$, and $\bd=\bg_m^{k}$, yields
\begin{align}
F(\bx^{k+1},\bB^{k})\leq F(\bx^{k},\bB^{k})-\left(\frac{1-\alpha_k}{\eta_k}-2L_F\right)\frac{\|\bx^{k+1}-\bx^k\|^2}{2}+\left(\frac{\alpha_k}{\eta_k}+2L_F\beta_k^2\right)\frac{\|\bx^k-\bx^{k-1}\|^2}{2}+\frac{1}{L_F}\|\nabla_{\bx}F(\bz^k,\bB^k)-\bg_m^k\|^2.
\end{align}
\end{bluequote}
\item (Lemma 5.8) I understand the connection between your algorithm and those presented in [62]. However, I think a brief proof would be welcome, particularly to explain why you do not obtain the same coefficients in the right-hand side of (5.20) as in the original paper.
\par\smallskip
\normalfont\upshape\textbf{Response:}
The submitted manuscript stated the StoRM descent estimate without displaying the coefficient construction or the resulting two component bounds.
We have added the coefficient identities and the proof reference in the StoRM complexity analysis \green{(Page 16, lines 451--458):}
\begin{bluequote}
In the next lemma, we establish an approximate sufficient descent property of $\Psi_r$, with the proof provided in Appendix~D.1.
\begin{align}
\begin{cases}
C_{\bB}=C_2-K_3=\frac{K_2-K_3-K_1}{2}>0,\\
C_{\bx}^{k}=\eta_k^{-1}-\widetilde L_F,\qquad \widetilde L_F=L_F+K_3+K_4.
\end{cases}
\end{align}
\end{bluequote}
\item Some proofs are placed in the main paper and others similar in the appendix (for example Lemma 5.4/Lemma 5.5). Is there a specific reason for this ?
\par\smallskip
\normalfont\upshape\textbf{Response:}
The submitted manuscript placed closely related one-step proofs in different locations without an explanatory pointer.
We have moved the proof to Appendix~C and added an explicit pointer from the main text \green{(Page 10, lines 160--165; Page 32, lines 951--956):}
\begin{bluequote}
The one-iteration progress on $\bx$ is estimated as follows, with the proof provided in Appendix~C.
\end{bluequote}
\item (page 10) ``the set $\mathrm{crit}(\mathscr{L})$... is nonempty'', I would place this fact in Assumption 5.1.
\item (Assumption 5.1) I would reformulate (ii) as ``$F$ is $L_F$-smooth ($L_F\in(0,+\infty)$) over $\mathbb{R}^q\times\mathbb{R}^{n\times d}$;''.
\par\smallskip
\normalfont\upshape\textbf{Response:}
The submitted manuscript stated nonemptiness of the critical set outside the assumptions and did not specify the domain and range of $L_F$ in Assumption~5.1(ii).
We have incorporated both points into Assumption~5.1 \green{(Page 11, lines 308--310):}
\begin{bluequote}
The critical set
\[
\mathrm{crit}\,\calL:=\{(\bx^*,\bB^*,\bLmbd^*)\in\RR^q\times\RR^{n\times d}\times\RR^{n\times d}:\bm{0}\in\partial\calL(\bx^*,\bB^*,\bLmbd^*)\}
\]
is nonempty.
\end{bluequote}
\item (page 13) For better consistency and since you state that $C_1,C_2,C_3,C_4$ are independent of $k$, you should introduce (5.14) setup before (5.12) and (5.13).
\par\smallskip
\normalfont\upshape\textbf{Response:}
The submitted manuscript introduced the $k$-independence of $C_1,C_2,C_3,C_4$ without stating the constant-parameter conditions that justify it.
We have made this dependence explicit in the Lyapunov-function definition \green{(Page 14, lines 398--399):}
\begin{bluequote}
where the positive coefficients $C_1,C_2,C_3,C_4$ (independent of $k$, since $\alpha_k\equiv\alpha$, $\beta_k\equiv\beta$, and $\tau$ is a constant; see (5.15) below) are chosen as:
\end{bluequote}
\item (Lemma C.1) I would replace $\nabla f(\bz)$ by a generic direction $\bd\in\mathbb{R}^q$, this doesn't change the proof. The reason is that, in the proof of Lemma 5.5, you apply (C.1) by considering $\bG^k$, which is not a gradient strictly speaking (but merely an oracle).
\par\smallskip
\normalfont\upshape\textbf{Response:}
The submitted Lemma~C.1 was written only for a gradient although the later proof applies it to a stochastic oracle direction.
We have restated Lemma~C.1 for an arbitrary direction $\bd\in\RR^q$ \green{(Page 33, lines 959--961):}
\begin{bluequote}
Lemma C.1. Let $f(\bx):\RR^q\to\RR$ be a continuously differentiable function with an $L$-Lipschitz continuous gradient. Then, for any $\bx,\by,\bd\in\RR^q$ and any $\bx^+\in\RR^q$ defined by $\bx^+=\by-\eta\bd$, $\eta>0$,
\end{bluequote}
\item (Lemma 5.5/proof) To better justify (C.3), simply say you apply Lemma C.1 to $F(.,\bB^k)$ considering $\bx=\bx_k,\by=\by_k,\eta=\eta_k,\bd=\bG^k$. The rest of the proof appears correct to me.
\par\smallskip
\normalfont\upshape\textbf{Response:}
The submitted proof did not state how the auxiliary lemma is instantiated for the stochastic $\bx$-update.
We have added the substitution $F(\cdot,\bB^k)$, $\bx=\bx^k$, $\by=\by^k$, $\eta=\eta_k$, and $\bd=\bg_m^k$ in the proof \green{(Page 33, lines 969--977):}
\begin{bluequote}
Then, applying the auxiliary lemma to $F(\cdot,\bB^k)$, which is $L_F$-smooth, with $\bx=\bx^k$, $\by=\by^k$, $\eta=\eta_k$, and $\bd=\bg_m^k$ yields the one-iteration estimate.
\end{bluequote}
\end{itemize}

\subsubsection{Notation and terminology}

\begin{itemize}\itshape
\item In some places, you use of letter $d_x$ to denote the space dimension relative to variable $x$ and sometimes letter $q$ instead. I would prefer the use of letter $q$ everywhere, $d_x$ notation suggests that the space dimension may change depending on the value of $x$ which is a bit confusing.
\par\smallskip
\normalfont\upshape\textbf{Response:}
The submitted manuscript used both $d_x$ and $q$ for the dimension of the parameter space.
We have replaced $d_x$ by $q$ consistently throughout the manuscript \green{(Page 7, lines 189--204):}
\begin{bluequote}
Deep hashing aims to learn network parameters $\bx\in\RR^{q}$ that generate binary codes preserving semantic similarity in the Hamming space.
\end{bluequote}
\item $h$ and $h^*$ functions are used for vectors and matrix objects. This is an acceptable abuse of notation in the context of your work provided that this is clearly stated, for example in your notation subsection.
\par\smallskip
\normalfont\upshape\textbf{Response:}
The submitted manuscript used $h$ and $h^*$ for scalar, vector, and matrix objects without stating the extension convention.
We have added the convention in the Notations paragraph \green{(Page 5, lines 81--82):}
\begin{bluequote}
For simplicity, we use the same symbols $h$ and $h^*$ for their scalar forms and their separable extensions to vectors and matrices whenever no confusion can arise. The Fenchel conjugate of $h$ is $h^*(\bLmbd):=\sup_{\bZ\in\RR^{n\times d}}\{\langle\bLmbd,\bZ\rangle-h(\bZ)\}$.
\end{bluequote}
\item Some `non-matrix' variables (for example, $\bD^k$, $\bG^k$) are written in capital letters, which does not strictly adhere to the conventions you set out in the notations subsection.
\par\smallskip
\normalfont\upshape\textbf{Response:}
The submitted manuscript used capital letters for stochastic direction vectors, conflicting with the notation convention for vectors.
We have renamed the two stochastic directions as $\bg_m^k$ and $\bg_s^k$ and updated their definitions \green{(Page 9, lines 260--267; Page 10, lines 348--360):}
\begin{bluequote}
\[
\bg_m^k:=\nabla_{\bx}F(\bz^k,\bB^k;\calJ_k)=\frac{1}{b_k}\sum_{i\in\calJ_k}\nabla_{\bx}F_i(\bz^k,\bb_i^k),
\qquad
\bx^{k+1}=\by^k-\eta_k\bg_m^k.
\]
\end{bluequote}
\item (page 5) Your definition of closeness for $f$ is equivalent to those of the lower semicontinuity (see for example chapter 1 of [1]. In the appendix, you mention the two terms (closeness in Def A.1, lower semicontinuity Prop A.2) . Because the second one, appears more popular to me, I would suggest to work with lower semicontinuity everywhere.
\par\smallskip
\normalfont\upshape\textbf{Response:}
The submitted manuscript used ``closed'' for an extended-real-valued function where ``lower semicontinuous'' is the intended terminology.
We have replaced the function terminology in the preliminary definitions and Appendix~A \green{(Page 6, lines 164--169; Page 30, lines 882--884):}
\begin{bluequote}
For a proper, lower semicontinuous, and convex function $f$, the proximal mapping of $f$ is defined as
\begin{align}
\operatorname{prox}_{\tau f}(\by):=\arg\min_{\bx\in\operatorname{dom}(f)}\left\{f(\bx)+\frac{1}{2\tau}\|\bx-\by\|^2\right\}.
\end{align}
\end{bluequote}
\item (page 28) See my first comment of preliminaries section about closeness vs lower-semicontinuity.
\par\smallskip
\normalfont\upshape\textbf{Response:}
The same ``closed'' versus ``lower semicontinuous'' inconsistency remained in Appendix~A.
We have revised Definition~A.1 and Proposition~A.2 to use ``lower semicontinuous'' consistently \green{(Page 30, lines 882--884):}
\begin{bluequote}
Let $f:\RR^{q_1}\to(-\infty,+\infty]$ and $g:\RR^{q_2}\to(-\infty,+\infty]$ be proper and lower semicontinuous.
\end{bluequote}
\item (Appendix B) I didn't really understand how you deduce, from the expression of $L_s$ that $L_F$ is $n$-independent. I would rather say that $L_F$ can be bounded by a constant independent of $n$. Following one of my previous comment, I don't really like the use of $\mathscr{O}$ notation speaking of constants.
\par\smallskip
\normalfont\upshape\textbf{Response:}
The submitted Appendix~B argument asserted $n$-independence of $L_F$ without displaying the constant bound.
We have introduced explicit $C_s$ and $C_F$ bounds and removed the informal $\mathscr{O}(1)$ statement \green{(Page 31, lines 902--915):}
\begin{bluequote}
Thus $C_s$ is finite and independent of $n$. Since $F=f+P$, its smoothness constant satisfies $L_F\leq L_f+C_s=:C_F$, where $C_F$ is finite and independent of $n$.
\end{bluequote}
\item Why do you keep $|\mathscr{I}_{k+1}|$ instead of using notation $b_{k+1}$?
\par\smallskip
\normalfont\upshape\textbf{Response:}
The submitted Appendix~D formulas used $|\mathcal J_{k+1}|$ although the algorithms define the batch size as $b_{k+1}$.
We have replaced the cardinality notation by $b_{k+1}$ throughout Appendix~D \green{(Page 39, lines 645--646):}
\begin{bluequote}
\[
\mathbb E[\|\be^{k+1}\|^2]\leq(1-\rho_k)^2\mathbb E[\|\be^k\|^2]+\frac{2\rho_k^2\sigma^2}{b_{k+1}}+\frac{4(1-\rho_k)^2L_F^2}{b_{k+1}}\mathbb E[\Delta_{\bx}^{k+1}+\Delta_{\bB}^{k+1}].
\]
\end{bluequote}
\item (page 7) Variable $\bm{b}_i$ denotes first $sgn(\bm{D}_i(\bx))$ and then, in subsection 3.2, corresponds to $\bm{D}_i(\bx)$, please clarify with two separated notations.
\par\smallskip
\normalfont\upshape\textbf{Response:}
The submitted manuscript used $\bb_i$ both for the binary code and for the continuous split variable.
We now reserve $\bb_i$ for the continuous auxiliary variable and write the binary code as $\bh_i=\operatorname{sgn}(\bD_i(\bx))$ \green{(Page 7, lines 197--204):}
\begin{bluequote}
The model is trained to produce the $i$-th continuous network output $\bD_i(\bx)$ with W-type regularization $h$, and binary codes are obtained post-training via $\bh_i=\operatorname{sgn}(\bD_i(\bx))\in\{-1,1\}^{d}$.
\end{bluequote}
\item (page 7) I would place definition in (4.1) in the subsection notations. Do not forget to mention the fact you use the same notations for both vectors and matrices (see 4th remark of ``General comments'' section)
\par\smallskip
\normalfont\upshape\textbf{Response:}
The submitted manuscript introduced the conjugate in the main algorithm section without fixing its matrix domain or stating the scalar/vector/matrix notation convention.
We have moved the definition to the Notations paragraph and corrected the domain to $\RR^{n\times d}$ \green{(Page 5, lines 81--82):}
\begin{bluequote}
The Fenchel conjugate of $h$ is $h^*(\bLmbd):=\sup_{\bZ\in\RR^{n\times d}}\{\langle\bLmbd,\bZ\rangle-h(\bZ)\}$ for $\bLmbd\in\RR^{n\times d}$.
\end{bluequote}
\item (page 6) Wouldn't it be ``$f(\bx)=\EE_{\epsilon}[F(\bx,\epsilon)]$'' instead ? Moreover, what function $F:(\bx,\epsilon)\mapsto F(\bx,\epsilon)$ is supposed to represent, is this a generic one ?
\par\smallskip
\normalfont\upshape\textbf{Response:} We thank the reviewer for identifying a notation conflict. The stochastic problem in Section~2 is a generic stochastic objective introduced only to explain the StoRM estimator.
We apologize for the confusion. To this end, the revision now in \green{Page 6, lines 1-3:} 
\begin{bluequote}
Consider a stochastic optimization problem with objective
\[
\phi(\bx)=\mathbb{E}_{\xi}[\phi(\bx;\xi)], \qquad \xi\in\Xi.
\]
A standard momentum-based stochastic first-order method constructs the stochastic direction
\begin{align*}
    \bg^{k} = \rho_{k-1}\nabla\phi(\bx^{k},\xi^{k}) +(1-\rho_{k-1})\bg^{k-1}, \qquad \rho_{k-1}\in[0,1],
\end{align*}
which is generally a biased estimator of $\nabla\phi(\bx^{k})$. The iterate is then updated according to
\[
\bx^{k+1}=\bx^{k}-\eta_k\bg^{k}.
\]
StoRM introduces the correction term
\[
(1-\rho_{k-1})\bigl(\nabla\phi(\bx^{k},\xi^{k})-\nabla\phi(\bx^{k-1},\xi^{k})\bigr)
\]
into the momentum estimator above, leading to the recursive estimator
\begin{align}
    \label{Eq: the generic StoRM estimator}
    \bg^{k}=\left\{\begin{array}{ll}
        \nabla\phi(\bx^{k},\xi^{k}), & \text{if } k=1, \\[1mm]
        \nabla\phi(\bx^{k},\xi^{k}) +(1-\rho_{k-1})\bigl(\bg^{k-1}-\nabla\phi(\bx^{k-1},\xi^{k})\bigr), & \text{if } k>1,
    \end{array}\right.
\end{align}
\end{bluequote}

\item (Prop A.2) In (iii), I would rather write $\hat{\partial}f_1(\overline{\bx}_1),\hat{\partial}f_2(\overline{\bx}_2)$ instead of $\hat{\partial}_{\bx_1}f_1(\overline{\bx}_1),\hat{\partial}_{\bx_2}f_2(\overline{\bx}_2)$
\par\smallskip
\par\smallskip
\normalfont\upshape\textbf{Response:}
The submitted Proposition~A.2(iii) used coordinate subscripts on the limiting subdifferentials that were not needed.
We have removed those subscripts in Appendix~A \green{(Page 30, lines 885--890):}
\begin{bluequote}
For any $\bx\in\RR^{q_1}$, we have $\hat{\partial}f(\bx)\subseteq\partial f(\bx)$; if $f$ is convex, then $\hat{\partial}f(\bx)=\partial f(\bx)$.
\end{bluequote}
\item (page 12) In the definition of $\Delta_B^k$, there is no need to specify the `$F$' for the norm as there is no ambiguity here.
\par\smallskip
\par\smallskip
\normalfont\upshape\textbf{Response:}
The submitted definition of $\Delta_{\bB}^k$ carried an unnecessary $F$ subscript on the norm.
We have removed the subscript in the complexity-analysis notation \green{(Page 10, lines 154--158):}
\begin{bluequote}
$\Delta_{\bx}^{k}=\frac{\|\bx^{k}-\bx^{k-1}\|^{2}}{2},\qquad \Delta_{\bB}^{k}=\frac{\|\bB^{k}-\bB^{k-1}\|^{2}}{2}$.
\end{bluequote}
\item (Lemma D.4) I would change notation $C_\delta^1$ by $C_\delta$ (you did this in the proof).
\par\smallskip
\par\smallskip
\normalfont\upshape\textbf{Response:}
The submitted Lemma~D.4 used the inconsistent symbol $C_\delta^1$ although the proof used $C_\delta$.
We have unified the notation as $C_\delta$ in Lemma~D.4 and its proof \green{(Page 40, lines 689--695):}
\begin{bluequote}
Moreover $\delta=\frac16$ and $c_\delta<36\left(\frac{\sqrt{30}-2}{13}\right)^2$ give $r=\tau m\leq\tau L_F<\frac{\sqrt{30}-2}{13}<\frac27$.
\end{bluequote}
\item (Remark 5.1) ``The smoothness constant $L_F$ remains $\mathscr{O}(1)$'', useless comment, a non-zero constant being always a $\mathscr{O}(1)$.
\par\smallskip
\par\smallskip
\normalfont\upshape\textbf{Response:}
The submitted Remark~5.1 used the uninformative phrase ``$L_F$ remains $\mathscr{O}(1)$''.
We have replaced it with an explicit bound independent of $n$ \green{(Page 11, lines 311--313):}
\begin{bluequote}
The smoothness constant $L_F$ is bounded by a constant independent of the sample size $n$.
\end{bluequote}
\end{itemize}

\par\smallskip
\normalfont\upshape\textbf{Response (notation summary).} Prompted by these comments, we re-examined the notation of the whole manuscript, appendices included, so that every symbol carries a single meaning. Table~\ref{tab:notation} lists the conventions of the revised manuscript; rows marked $^\star$ concern clashes that we found ourselves.

\begin{table}[!ht]
\centering\footnotesize\color{blue}
\renewcommand{\arraystretch}{1.2}
\caption{Notation conventions in the revised manuscript.}\label{tab:notation}
\begin{tabular}{@{}p{2.5cm}p{7.6cm}p{7.6cm}@{}}
\hline
\textbf{Symbol} & \textbf{Meaning in the revision} & \textbf{Change from the submitted version}\\
\hline
\multicolumn{3}{@{}l}{\textit{Dimensions and indices}}\\
$q$ & dimension of the network parameter $\bx\in\mathbb{R}^q$ & $d_x$ replaced by $q$ throughout\\
$\ba^{\mathrm{qry}}$, $\bh^{\mathrm{qry}}$, $R^{\mathrm{qry}}$ & an unseen query, its binary code, and its number of relevant items (Appendix~E) & previously $\ba_q$, $\bb_q$, $R_q$, and $\bx_q$ in the mAP definition; $q$ is now reserved for the parameter dimension$^\star$\\
\hline
\multicolumn{3}{@{}l}{\textit{Network outputs, auxiliary variables, and codes}}\\
$\bC_i(\bx)$, $\bD_i(\bx)$, $\mathcal{D}(\bx)$ & $\bC_i(\bx)=\bC(\bx;\ba_i)$ is the output before activation; $\bD_i(\bx)\in\mathbb{R}^d$ is the continuous output of sample $i$; $\mathcal{D}(\bx)=(\bD_1(\bx),\ldots,\bD_n(\bx))$ & $\bu_i$ of Appendix~E is written $\bD_i(\bx)$, since both denote the same object; the domain of the mapping in Section~2, formerly $\mathcal{D}$, is now $\Omega$, so that $\mathcal{D}$ always denotes the stacked map$^\star$\\
$\bb_i$, $\bB$ & continuous auxiliary variable $\bb_i\in\mathbb{R}^d$, the $i$-th row of $\bB=[\bb_1^\top;\ldots;\bb_n^\top]\in\mathbb{R}^{n\times d}$, tied to $\bD_i(\bx)$ by the penalty $P$ & $\bb_i$ no longer denotes the binary code (Section~3, Appendix~E); $\bB$ was written $[\bb_1,\ldots,\bb_n]$, which suggests columns$^\star$\\
$\bh_i$ & binary hash code $\bh_i=\operatorname{sgn}(\bD_i(\bx))\in\{-1,1\}^d$ & replaces $\bb_i$ as the binary code in Section~3 and Appendix~E, where $\bh_i$ already denotes the code in the likelihood\\
\hline
\multicolumn{3}{@{}l}{\textit{Functions}}\\
$f$, $P$, $F$ & $f(\bx)=n^{-1}\sum_i f_i(\bx)$; $P(\bx,\bB)=\frac{\gamma}{2n}\sum_i\|\bD_i(\bx)-\bb_i\|^2$; $F=f+P$ & the undefined function $G$ in the proof of Lemma~5.2 is removed; the optimality conditions read $\nabla_{\bx}F(\bx^*,\bB^*)=\bm{0}$ and $\nabla_{\bB}F(\bx^*,\bB^*)+\bLmbd^*=\bm{0}$\\
$h$, $h^*$ & $h(t)=\lambda\,\bigl||t|-1\bigr|$ for scalars, extended separably to vectors and matrices; $h^*(\bLmbd)=\sup_{\bZ\in\mathbb{R}^{n\times d}}\{\langle\bLmbd,\bZ\rangle-h(\bZ)\}$ & definition (4.1) moved to the Notations paragraph; its domain, formerly $\mathbb{R}^n$, is corrected to $\mathbb{R}^{n\times d}$$^\star$\\
$\mathcal{L}$, $\ell_{\mathcal{S}}$ & $\mathcal{L}(\bx,\bB,\bLmbd)=F(\bx,\bB)+\langle\bB,\bLmbd\rangle-h^*(\bLmbd)$ is the Lagrangian; $\ell_{\mathcal{S}}(\bx)$ is the pairwise loss of Appendix~E & the pairwise loss, formerly $\mathcal{L}_S$, is renamed$^\star$\\
\hline
\multicolumn{3}{@{}l}{\textit{Algorithm and analysis}}\\
$\bg_m^k$, $\bg_s^k$, $\bg^k$ & stochastic directions of DualHash-StoM and DualHash-StoRM, and the generic momentum estimator of Section~2 & replace $\bG^k$, $\bD^k$, and the generic $\bd^k$ of Section~2\\
$\bd^k$ & $\bd^k=(\bd_{\bx}^k,\bd_{\bB}^k,\bd_{\bLmbd}^k)\in\partial\mathcal{L}(\bx^k,\bB^k,\bLmbd^k)$ with $\bd_{\bx}^k=\nabla_{\bx}F(\bx^k,\bB^k)$, $\bd_{\bB}^k=\nabla_{\bB}F(\bx^k,\bB^k)+\bLmbd^k$, $\bd_{\bLmbd}^k=\bB^k-\mathcal{G}_{\bLmbd}^k$, $\mathcal{G}_{\bLmbd}^k\in\partial h^*(\bLmbd^k)$; used only in the stationarity measure & the definition is now numbered and $\bd^k$ keeps this meaning only, so $\bd^k$ and the update directions $\bg^k$ are different objects\\
$\alpha_k$, $\beta_k$, $\kappa$ & $\alpha_k,\beta_k$ are the momentum parameters of DualHash-StoM; $\kappa$ is the scale of the sigmoid $\operatorname{sig}_\kappa(t)=(1+e^{-\kappa t})^{-1}$ in Appendix~E & the sigmoid scale, formerly $\alpha$, is renamed $\kappa$$^\star$\\
$\sigma^2$, $\Phi$ & $\sigma^2$ is the variance bound in Assumption~5.1(iii); $\Phi(\bx_1,\bx_2)=f_1(\bx_1)+f_2(\bx_2)$ in Proposition~A.2 & $\sigma$ also denoted the sigmoid and the sum in Proposition~A.2; they are now $\operatorname{sig}_\kappa$ and $\Phi$$^\star$\\
\hline
\end{tabular}
\end{table}


\subsubsection{References}

\begin{itemize}\itshape
\item (page 2) ``using subgradients may result in slow convergence and oscillation.'' add references to illustrate such a fact.
\par\smallskip
% 中文注释：怕得罪次微分那帮人。
\normalfont\upshape\textbf{Response:} We thank the reviewer for this suggestion. We have added references documenting the slower convergence rates of stochastic subgradient methods for nonsmooth optimization and the oscillatory dynamics of subgradient methods, namely Shamir and Zhang (2013) and Bolte, Pauwels, and R{\'\i}os-Zertuche (2024).
\item (page 2) ``inverse imaging problems'' I would also add additional references, in the vein of [3] et al.
\item (page 6) Add reference [2] for descent lemma.
\par\smallskip
\normalfont\upshape\textbf{Response:} We thank the reviewer for this suggestion. We have added the citation to Bertsekas [2] at the first occurrence of the descent lemma on page 6.
\item (page 17) including pointwise [...] methods'' add references.
\par\smallskip
\normalfont\upshape\textbf{Response:} We thank the reviewer for this suggestion. We have added representative references for pointwise, pairwise, and tripletwise supervised hashing methods: SDH, DSH, and DTSH, respectively.
\item (Prop A.3) I agree with the results of (i), (ii) and (iii) but the corresponding lemma/article (Lemma 5.2 [45]) doesn't not appear to be the correct reference.
\par\smallskip
\normalfont\upshape\textbf{Response:} We thank the reviewer for identifying this reference error. The relations stated in Proposition A.3 correspond to Proposition 4.1 of Pock and Sabach's iPALM paper, rather than Lemma 5.2. We have corrected the citation accordingly. The form used in our proposition writes the relations directly in terms of squared norms, whereas the original paper uses $\Delta_i^k=\frac{1}{2}\|\bx_i^k-\bx_i^{k-1}\|^2$.
\end{itemize}

\subsubsection{Wording}

\begin{itemize}\itshape
\item (page 1) I would reformulate as ``Mathematically, this amounts \textbf{to carry out} a projection onto [...]. \textbf{Such a process can be interpreted as} a nonconvex discrete optimization problem[...], where the date fidelity criterion \textbf{is based on} pairwise similarity...''
\item (page 3) ``[...] deep hashing; directly extending their results to our setting is nontrivial.'' I would cut the sentence in two, for example : ``[...] deep-hashing. A direct extension of their results to our setting is far from being trivial. ''
\item (page 11) ``The presence [...] renders the Lagrangian (4.2) non-descent'', I understand what you're saying, but the wording is still a bit vague. I would reformulate as ``The presence of dual variables $\Lambda$ makes it difficult to establish a descent property for the Lagrangian''.
\item (page 13) ``To handle the oscillation of (4.2). I would reformulate as ``To handle the oscillations of the Lagrangian function $\mathscr{L}$''.
\end{itemize}

\par\smallskip
\normalfont\upshape\textbf{Response:} We thank the reviewer for these suggestions. All four passages have been revised as suggested; for the second one, the revised sentence also says why the extension to the stochastic setting is far from trivial.

\subsubsection{Experiments}

\begin{itemize}\itshape
\item \red{(page 18) Is there a particular reason why you are using a step size of $0.1\tau$ rather than $\tau$ when updating $\Lambda_k$?}
\item \red{It would have been interesting to also display, in Tables 6.2 and 6.3, the performances of DualHash-StoRM.}
\item \red{I think you could try increasing the size of a few figures (especially 6.4).}
\end{itemize}

\subsection{Typographical and formatting issues}

\par\smallskip
\textbf{General response:} We thank the reviewer for the careful reading. We have corrected the typographical, formatting, notation, and cross-reference issues listed below and proofread the manuscript to identify and fix similar errors elsewhere. Additional responses are provided below for the comments that require clarification of the revision or its calculation.

\begin{itemize}\itshape
\item (page 1) ``$W$-type regularization such as $z\mapsto \lvert |z|-1\rvert$''.
\item (page 1) (Typo ?) ``up to $512\times$''
\par\smallskip
\normalfont\upshape\textbf{Response:} We thank the reviewer for pointing out that this statement was unclear. The figure $512\times$ was intended as an illustrative example of the storage efficiency of hashing. Specifically, representing a 1024-dimensional \texttt{float32} feature using a packed 64-bit binary code reduces the per-sample representation storage by a factor of $(1024\times32)/64=512$. This factor depends on the feature dimension, numerical precision, and code length. We have replaced ``up to $512\times$'' with an explicit example and added a footnote giving the calculation and clarifying that it excludes model parameters and index overhead.
\item (page 1) ``such as $z\mapsto \lvert |z|-1\rvert$''.
\item (page 2-3) ``reformulate the problem as \textbf{one of constrainted optimization}''.
\item \red{(page 3) ``the challenging composite term $h\circ \mathbf{D}$''}
\item (page 4/Related work) ``W-type regularizer such as $z\mapsto \lVert|z|-1\rVert_1$'', ``employ $z\mapsto \lVert|z|-1\rVert_2^2$ as a smooth surrogate for $z\mapsto \lVert|z|-1\rVert_1$''.
\item (page 5) ``and \sout{the} stochastic majorization-minimization subspace methods''.
\item (page 6) ``A \uline{differentiable} function $f:\mathbb{R}^q\to\mathbb{R}$ is said to be $L_f$-smooth''.
\item (page 13) ``descent property of $\underline{\{}\Psi_m(\mathscr{Q}^k)\underline{\}}$''.
\item (page 13) ``\uline{are} established in Lemma C.2''.
\item ``From (5.25), to \uline{achieve}''.
\item (page 17) ``regularization term $h(\bz)=\lambda\lVert|\bz|-1\rVert_1$ (the first z was not in bold).
\item (page 17) ``are provided in the Supplementary Materials \uline{E}''.
\item (page 17) ``to observe convergence behaviour\uline{s}''.
\item (page 22) ``helps \uline{to} produce high mAP''.
\item ``including image \uline{reconstruction}''.
\item ``and binary neural networks'' (no comma).
\item (page 28) In the equality just before Prop A.2, one of the $x$ is not in bold.
\item (Lemma C.3/proof) ``is deduced from the $\underline{L_F}$-smoothness of $\underline{F(.,\bB^k)}$''.
\par To justify (C.15), I would mention inequality $(a+b)^2\le 2a^2+2b^2$ just after the $F$ smoothness (and not after the inequality).
\par ``\sout{yield} conduct to (C.9).''
\par To justify (C.16), I would have mentioned Lemma 5.3 directly instead of (5.6).
\end{itemize}

\subsection{Clarification and possible misunderstanding}

\begin{itemize}\itshape
\item (page 2) ``its composition with neural network outputs $h(\mathbf{D}i(x))$'' sounds a bit unclear to me.
\par\smallskip
\normalfont\upshape\textbf{Response:} We thank the reviewer for pointing out this ambiguity. Here $\mathbf{D}_i(\bx)$ denotes the continuous network output for the $i$-th sample, and $h\circ\mathbf{D}_i$ denotes the composition $\bx\mapsto h(\mathbf{D}_i(\bx))$. We have revised the sentence to make this notation explicit.
\item (page 2) ``when $h$ is component-wise nondecreasing'', to verify but after quickly checked [29], I think $h$ also need to be convex. This would corroborate your next sentence ``when $h$ itself is nonconvex...''.
\par\smallskip
\normalfont\upshape\textbf{Response:} We agree that the previous wording conflated two different sufficient conditions. We now state the assumptions separately: the first cited condition uses convex $h$ and smooth $\bD$, while the cited monotonicity-based result assumes a weakly convex and component-wise nondecreasing outer function together with a weakly convex inner mapping. We also clarify that our nonconvex W-type regularizer is outside these sufficient-condition results.
\item (page 8) What do you mean by ``inexpensive gradient'' ?
\par\smallskip
\normalfont\upshape\textbf{Response:} We have replaced this informal phrase with an explicit explanation. Since $f$ does not depend on $\bB$, the gradient is $\nabla_{\bB}\calL(\bx,\bB,\bLmbd)=\frac{\gamma}{n}(\bB-\mathcal{D}(\bx))+\bLmbd$. It is computed element-wise once the network outputs are available and does not require backpropagation with respect to $\bx$.
\end{itemize}



\endgroup
\section{Lower bounds for the Lyapunov functions}\label{sec:lowerbounds}
% CHECK before moving into the manuscript: (a) DONE (retuned constants + explicit eta-conditions in the Remark): conditions (P1),(P2),(P1'),(P2') are smallness conditions on delta and eta,
%      verify they are compatible with tau < (sqrt(30)-2)/(13 L_F) of Lemma C.2 and the eta of (5.14);
%      (b) DONE: Lambda^1 in [-lambda,lambda]^{n x d} is now part of the initialization of Algorithms 4.1 and 4.2 (methodology.tex);
%      (c) StoRM part must be aligned with the repaired schedule of Theorem 5.3 (see the appeal letter, Comments 2--3).
\begingroup\color{blue}
Throughout this section, $m:=\gamma/n$, $M:=\lambda\sqrt{nd}$, $\mathcal{D}(\bx):=(\bD_1(\bx),\ldots,\bD_n(\bx))$, and we use the convention $\bx^0:=\bx^1$. We write $\bm{r}_k:=\bB^k-\mathcal{D}(\bx^k)$, $\bm{p}_k:=\bx^k-\bx^{k-1}$, and $\bm{u}_k:=\bx^{k+1}-\bx^k$. We use the output condition of Appendix~B: $\|\mathcal{D}(\bx)\|_F\le C_{\mathcal{D}}$ and $\mathcal{D}$ is $L_{\mathcal{D}}$-Lipschitz. The constants $C_1,\dots,C_4$ are those of the Lyapunov functions $\Psi_m$ in~(5.12) and $\Psi_r$ in~(5.21) of the submitted manuscript, and $\alpha,\beta,\eta,\tau$ are the parameters of Theorem~5.2. All expectations are assumed finite.

\begin{proposition}\label{prop:lyapunov-lower-bound}
Let Assumption~5.1 hold and let $\bLmbd^1\in[-\lambda,\lambda]^{n\times d}$, which is now required in the initialization of Algorithms~4.1 and~4.2. Set
\[
c_0:=\tfrac{4M^2}{m}+MC_{\mathcal{D}}+\lambda nd,\qquad
A_B:=\tfrac{3}{2}C_1\tau^2m^2,\qquad
H_m:=\tfrac{3}{2}C_1\tau^2m^2L_{\mathcal{D}}^2+\tfrac{1}{2}C_3 .
\]
\begin{enumerate}[label=\textup{(\roman*)},nosep]
\item \textup{(DualHash-StoM)} Let $\eta_k\equiv\eta/L_F$, $\alpha_k\equiv\alpha$, $\beta_k\equiv\beta$, and put $\Theta:=12H_m\eta^2/L_F$, $Q:=3\alpha^2+6\eta^2\beta^2$. If
\begin{align}
\tfrac{m}{2}(1-\Theta)-\tfrac{m}{16}-A_B&\ \ge\ \tfrac{m}{8},\tag{P1}\label{eq:P1}\\
C_4&\ \ge\ 2H_mQ,\tag{P2}\label{eq:P2}
\end{align}
then for every $k\ge1$,
\[
\mathbb{E}[\Psi_m(\calQ^k)]\ \ge\ \underline{\ell}+\tfrac{\mu}{8}\mathbb{E}\|\bx^k\|^2+\tfrac{m}{8}\mathbb{E}\|\bm{r}_k\|_F^2-c_1\ \ge\ \Psi_m^*:=\underline{\ell}-c_1,
\]
where $c_1:=c_0+\tfrac{3}{2}C_1\tau^2M^2+3H_m\eta^2\sigma^2/L_F^2$ does not depend on $T$.
\item \textup{(DualHash-StoRM)} Let $\eta_k\equiv\eta_T$, let $\be^k:=\bg_s^k-\nabla_{\bx}F(\bx^k,\bB^k)$, and put $\Gamma:=2H_m\eta_T^2$ and $W_k:=\mathbb{E}[\Psi_r(\calQ^k)]+\frac{\eta_T}{\rho_T}\mathbb{E}\|\be^k\|^2$. If
\begin{align}
\Gamma&\ \le\ \tfrac{\eta_T}{\rho_T},\tag{P1$'$}\label{eq:P1s}\\
\tfrac{m}{2}(1-2L_F\Gamma)-\tfrac{m}{16}-A_B&\ \ge\ \tfrac{m}{8},\tag{P2$'$}\label{eq:P2s}
\end{align}
then for every $k\ge1$,
\[
W_k\ \ge\ \underline{\ell}+\tfrac{\mu}{8}\mathbb{E}\|\bx^k\|^2+\tfrac{m}{8}\mathbb{E}\|\bm{r}_k\|_F^2-c_1'\ \ge\ \Psi_r^*:=\underline{\ell}-c_1',
\qquad c_1':=c_0+\tfrac{3}{2}C_1\tau^2M^2 .
\]
\end{enumerate}
\end{proposition}

\medskip\noindent\textit{Remark (conditions for the schedule of Theorem~5.3).}
For $\eta_T=\eta/(\tilde L_FT^{1/3})$ and $\rho_T=8\rho\eta^2/T^{2/3}$, $\Gamma$ decreases and $\eta_T/\rho_T$ increases with $T$, so (P1$'$) and (P2$'$) hold for every $T\ge1$ as soon as they hold at $T=1$, i.e.,
\[
16\rho H_m\eta^3\ \le\ \tilde L_F
\qquad\text{and}\qquad
\frac{2L_FH_m\eta^2}{\tilde L_F^2}+\frac{A_B}{m}\ \le\ \frac5{16}.
\]
With $\delta=\frac16$ we have $C_1=\frac{1}{2\tau}$ and $C_3=\frac{\tau L_F^2}{2}$, hence $A_B/m=\frac34\tau m$ and $H_m=\frac34\tau m^2L_{\mathcal{D}}^2+\frac14\tau L_F^2$. Since $m\le L_F$ and $\tau L_F<\frac{\sqrt{30}-2}{13}$, we get $A_B/m<0.21$, so the second condition holds whenever $2L_FH_m\eta^2\le\tilde L_F^2/10$. Both are smallness conditions on $\eta$; they are added to the constants of Theorem~5.3 when the proposition is moved into the manuscript.

\begin{proof}
\emph{Common estimates.} Since $f_i=\ell_{s,i}+\frac{\mu}{2}\|\cdot\|^2$ and $P(\bx,\bB)=\frac m2\|\bB-\mathcal{D}(\bx)\|_F^2$, Assumption~5.1(i) gives
\begin{equation}\label{eq:F-lower}
F(\bx,\bB)-\underline{\ell}\ \ge\ \tfrac{\mu}{2}\|\bx\|^2+\tfrac m2\|\bB-\mathcal{D}(\bx)\|_F^2\ \ge\ 0 .
\end{equation}
Because $F$ is $L_F$-smooth and bounded below by $\underline{\ell}$, $\|\nabla F(\bx,\bB)\|^2\le2L_F\bigl(F(\bx,\bB)-\underline{\ell}\bigr)$ for all $(\bx,\bB)$. The function $h$ is separable with $h^*(s)=|s|$ on $[-\lambda,\lambda]$ for each entry and $h^*=+\infty$ elsewhere, so $\operatorname{dom}h^*=[-\lambda,\lambda]^{n\times d}$ and $0\le h^*(\bLmbd)\le\lambda nd$ there. By the dual update, $\bLmbd^k\in[-\lambda,\lambda]^{n\times d}$ for all $k\ge1$, hence $\|\bLmbd^k\|_F\le M$.

\emph{Step 1: the Lagrangian.} As $\|\bB^k\|_F\le\|\bm r_k\|_F+C_{\mathcal{D}}$, Young's inequality gives $\langle\bB^k,\bLmbd^k\rangle\ge-M\|\bm r_k\|_F-MC_{\mathcal{D}}\ge-\frac m{16}\|\bm r_k\|_F^2-\frac{4M^2}{m}-MC_{\mathcal{D}}$. With $h^*\le\lambda nd$,
\begin{equation}\label{eq:L-lower}
\calL(\bx^k,\bB^k,\bLmbd^k)\ \ge\ F(\bx^k,\bB^k)-\tfrac m{16}\|\bm r_k\|_F^2-c_0 .
\end{equation}

\emph{Step 2: the forward $\bB$-increment.} The $\bB$-update and $\nabla_{\bB}F(\bx,\bB)=m(\bB-\mathcal{D}(\bx))$ give $\Delta_B^{k+1}=\frac{\tau^2}{2}\|m(\bB^k-\mathcal{D}(\bx^{k+1}))+\bLmbd^k\|_F^2$. Since $\bB^k-\mathcal{D}(\bx^{k+1})=\bm r_k+\mathcal{D}(\bx^k)-\mathcal{D}(\bx^{k+1})$ and $\mathcal{D}$ is $L_{\mathcal{D}}$-Lipschitz,
\[
\Delta_B^{k+1}\ \le\ \tfrac{3\tau^2}{2}\bigl(m^2\|\bm r_k\|_F^2+m^2L_{\mathcal{D}}^2\|\bm u_k\|^2+M^2\bigr).
\]
With $\Delta_x^{k+1}=\frac12\|\bm u_k\|^2$ this yields
\begin{equation}\label{eq:B-inc}
-C_1\Delta_B^{k+1}-C_3\Delta_x^{k+1}\ \ge\ -A_B\|\bm r_k\|_F^2-H_m\|\bm u_k\|^2-\tfrac32C_1\tau^2M^2 .
\end{equation}

\emph{Proof of \textup{(i)}.} Let $\bz^k=\bx^k+\beta\bm p_k$ and $\be^k:=\bg_m^k-\nabla_{\bx}F(\bz^k,\bB^k)$, so that $\mathbb{E}\|\be^k\|^2\le\sigma^2/b_k\le\sigma^2$. The $\bx$-update gives $\bm u_k=\alpha\bm p_k-\eta_k\bigl(\nabla_{\bx}F(\bz^k,\bB^k)+\be^k\bigr)$. Using $\|\nabla_{\bx}F(\bz^k,\bB^k)\|^2\le2\|\nabla_{\bx}F(\bx^k,\bB^k)\|^2+2L_F^2\beta^2\|\bm p_k\|^2$ and the gradient bound above,
\[
\mathbb{E}\|\bm u_k\|^2\ \le\ 12\eta_k^2L_F\,\mathbb{E}\bigl[F(\bx^k,\bB^k)-\underline{\ell}\bigr]+Q\,\mathbb{E}\|\bm p_k\|^2+3\eta_k^2\sigma^2 .
\]
Insert this and \eqref{eq:L-lower}--\eqref{eq:B-inc} into $\mathbb{E}\Psi_m(\calQ^k)=\mathbb{E}\calL-C_1\mathbb{E}\Delta_B^{k+1}+C_2\mathbb{E}\Delta_B^k-C_3\mathbb{E}\Delta_x^{k+1}+C_4\mathbb{E}\Delta_x^k$ and use $\|\bm p_k\|^2=2\Delta_x^k$:
\[
\begin{aligned}
\mathbb{E}\Psi_m(\calQ^k)\ \ge\ &\underline{\ell}+(1-\Theta)\mathbb{E}\bigl[F(\bx^k,\bB^k)-\underline{\ell}\bigr]-\bigl(\tfrac m{16}+A_B\bigr)\mathbb{E}\|\bm r_k\|_F^2\\
&+C_2\mathbb{E}\Delta_B^k+(C_4-2H_mQ)\mathbb{E}\Delta_x^k-c_0-\tfrac32C_1\tau^2M^2-3H_m\eta_k^2\sigma^2 .
\end{aligned}
\]
By \eqref{eq:P1}, $\Theta\le\frac58$, so \eqref{eq:F-lower} gives $(1-\Theta)(F-\underline{\ell})\ge\frac{3\mu}{16}\|\bx^k\|^2+(1-\Theta)\frac m2\|\bm r_k\|_F^2$, and \eqref{eq:P1} turns the coefficient of $\mathbb{E}\|\bm r_k\|_F^2$ into at least $\frac{m}{8}$ (and $\frac{3\mu}{16}\ge\frac{\mu}{8}$). The terms $C_2\mathbb{E}\Delta_B^k$ and $(C_4-2H_mQ)\mathbb{E}\Delta_x^k$ are nonnegative by \eqref{eq:P2}. Finally $3H_m\eta_k^2\sigma^2=3H_m\eta^2\sigma^2/L_F^2$ is independent of $T$.

\emph{Proof of \textup{(ii)}.} Here $\bm u_k=-\eta_T\bg_s^k=-\eta_T\bigl(\nabla_{\bx}F(\bx^k,\bB^k)+\be^k\bigr)$, so $\mathbb{E}\|\bm u_k\|^2\le2\eta_T^2\bigl(\mathbb{E}\|\nabla_{\bx}F(\bx^k,\bB^k)\|^2+\mathbb{E}\|\be^k\|^2\bigr)$. With \eqref{eq:L-lower}--\eqref{eq:B-inc}, $\Gamma=2H_m\eta_T^2$, and $\Psi_r(\calQ^k)=\calL-C_1\Delta_B^{k+1}+C_2\Delta_B^k-C_3\Delta_x^{k+1}$,
\[
\begin{aligned}
W_k\ \ge\ &\mathbb{E}F(\bx^k,\bB^k)-\bigl(\tfrac m{16}+A_B\bigr)\mathbb{E}\|\bm r_k\|_F^2-\Gamma\,\mathbb{E}\|\nabla_{\bx}F(\bx^k,\bB^k)\|^2\\
&+\bigl(\tfrac{\eta_T}{\rho_T}-\Gamma\bigr)\mathbb{E}\|\be^k\|^2+C_2\mathbb{E}\Delta_B^k-c_1' .
\end{aligned}
\]
The gradient bound gives $\Gamma\|\nabla_{\bx}F\|^2\le2L_F\Gamma(F-\underline{\ell})$, and \eqref{eq:P2s} implies $2L_F\Gamma\le\frac58$. Hence, as in (i), $F-2L_F\Gamma(F-\underline{\ell})-(\frac m{16}+A_B)\|\bm r_k\|_F^2\ge\underline{\ell}+\frac{3\mu}{16}\|\bx^k\|^2+\frac{m}{8}\|\bm r_k\|_F^2$. The error term has coefficient $\frac{\eta_T}{\rho_T}-\Gamma\ge0$ by \eqref{eq:P1s}, and $C_2\mathbb{E}\Delta_B^k\ge0$.
\end{proof}
\endgroup

\section*{References}

\begin{enumerate}
\item H. H. Bauschke and P. L. Combettes. \emph{Convex Analysis and Monotone Operator Theory in Hilbert Spaces}. CMS Books in Mathematics. Springer, Cham, 2 edition, 2017.

\item D. P. Bertsekas. Nonlinear programming. \emph{Journal of the Operational Research Society}, 48(3):334--334, 1997.

\item A. Repetti and Y. Wiaux. Variable metric forward-backward algorithm for composite minimization problems. \emph{SIAM Journal on Optimization}, 31(2):1215--1241, 2021.
\end{enumerate}

\end{document}
